\subsection{The relative normalizer comparison under the manuscript's exactness hypothesis}

The relative tangential construction in Deligne's treatment supplies
the comparison left open in the local calculation. We state its
hypotheses explicitly, rather than asserting it for an arbitrary
fibration or in mixed characteristic.

\begin{theorem}[F149-RELATIVE]\label{thm:archive-relative}
Let $S$ be a connected smooth scheme of finite type over a field of
characteristic zero. Let $X\to S$ be a smooth proper relative curve,
let $D\subset X$ be finite etale over $S$, and suppose that the
geometric fibres of $U=X\setminus D$ are hyperbolic. Fix a section
$g:S\to D$ and assume the exact homotopy sequence required in scan~79:
\[
 1\longrightarrow\pi_1(U_{\bar s})
 \longrightarrow\pi_1(U)\longrightarrow\pi_1(S)\longrightarrow1.
                                                               \tag{R149.1}
\]
Let $\mathcal L=J/J^2$ be the conormal line of $g$, let
$N=\mathcal L^\vee$ be its normal line, and let $N^\times$
be the complement of the zero section in its total space.
Choose the compatible tangential base points, and write
$I\simeq\widehat{\mathbb Z}(1)$ for geometric cusp inertia.
Then tangential specialization identifies extensions
\[
 \pi_1(N^\times)\ \simeq\ N_{\pi_1(U)}(I).
\]
\[
 1\longrightarrow I\longrightarrow\pi_1(N^\times)
 \longrightarrow\pi_1(S)\longrightarrow1.                       \tag{R149.2}
\]
It also identifies the tame formal-tube extension with these groups.
At each invertible level $n$, the image in etale cohomology of the
extension pushed out to $\mu_n$ is the Kummer Chern class of the
conormal line, up to the normal/conormal and inertia-orientation sign
already specified in F149-CHERN-LOCAL.

If $S$ is an etale $K(\pi,1)$ and the torsion subgroup of
$\operatorname{Pic}(S)$ has bounded exponent, this is the corresponding equality in
$H^2_{\rm cont}(S,\widehat{\mathbb Z}(1))$ under the usual comparison.
If $\operatorname{Pic}(S)$ is finitely generated, the extension in
(R149.2) splits if and only if $\mathcal L$ is trivial. A global
conormal basis supplies the resulting geometric splitting, and changes
of basis act by the uncompleted unit Kummer map of F149-SPLIT.
\end{theorem}

\begin{proof}
Deligne constructs the relative tangential specialization functor
from tame covers of $U$ to covers of the punctured relative tangent
line in \S15.19, printed p.~239 of \cite{DeligneTangential1989}.
His \S15.26, p.~242, explicitly applies the formal-coordinate
description to that relative situation. The construction is
functorial over $S$ and restricts, at a geometric point of $S$, to
the usual inertia map of the punctured tangent line into the curve.
Here all finite etale covers of $U$ are tame along the boundary:
$X$ is regular, the boundary is a relative normal-crossings divisor,
and every residue characteristic is zero. The local description and
descent are the generalized Abhyankar lemma in \emph{SGA~1}, XIII,
5.2--5.6 \cite{SGA1}. Thus the specialization functor is an exact
functor of the relevant Galois categories. It induces
\[
 \phi:\pi_1(N^\times)\longrightarrow\pi_1(U)
\]
over $G=\pi_1(S)$, with restriction to geometric fibre inertia
equal to the cusp inclusion $I\hookrightarrow\pi_1(U_{\bar s})$.
The formal-tube category here is the category of covers tame along
the removed section, as in the cited relative construction; no
category of wild covers is inserted.

The line bundle compactifies as a $\mathbb P^1$-bundle with the
disjoint zero and infinity sections. Hence the right-exact homotopy
sequence of \emph{SGA~1}, XIII, Proposition~4.1 and the second
example in~4.4 (reedition pp.~306--307 and~309), gives
\[
 I\longrightarrow\pi_1(N^\times)\longrightarrow G
 \longrightarrow1.                                           \tag{R149.3}
\]
The set of allowed primes is all primes in characteristic zero,
so the source's relative prime-set quotient is the full group here.
The composite of the first arrow with $\phi$ is injective: cusp
inertia injects in the geometric hyperbolic fundamental group, and
that group injects by (R149.1). Thus the first arrow of (R149.3)
is injective as well. This proves the exact sequence in (R149.2).
In particular no additional unproved assertion about a boundary map
from $\pi_2(S)$ is used; the exact sequence (R149.1) is precisely
the retained hypothesis.

Geometric cusp inertia is self-normalizing in
$\pi_1(U_{\bar s})$, by the geometric form of Stix's
Proposition~16, p.~8 \cite{StixCuspidal2012}.
The global boundary section preserves its geometric conjugacy class.
For any element of $G$, lift it to $e\in\pi_1(U)$ using
(R149.1). There is then $a\in\pi_1(U_{\bar s})$ such that
$eIe^{-1}=aIa^{-1}$, and $a^{-1}e$ belongs to the normalizer.
Consequently that normalizer surjects onto $G$, and its intersection
with the geometric kernel is exactly $I$. It therefore has an exact
sequence with the same kernel and quotient as (R149.3).

The subgroup $I$ is normal in $\pi_1(N^\times)$, so $\phi$
takes values in this normalizer. The resulting map of extensions
induces the identity on $I$ and on $G$ and is an isomorphism.
Indeed its kernel lies in $I$, where it is trivial. To lift an
element of the target, first lift its class in $G$ and then correct
by an element of the common kernel $I$. This proves surjectivity;
the continuous bijection is a homeomorphism because the groups
are compact and Hausdorff. Deligne's relative formal-coordinate
equivalence identifies the tame tube group with the same source,
proving the asserted tube comparison.

For the obstruction, choose local conormal frames, represented by
parameters $t_a$, with leading transition units $u_{ab}$. At level
$n$ the fibre cover is obtained by adjoining an $n$th root of the
parameter. The choices of roots of $u_{ab}$ identify these covers
on overlaps; their failure to satisfy descent on triples is the
same $\mu_n$-valued boundary cocycle as in F149-CHERN-LOCAL.
More explicitly, the gerbe of $\mu_n$-torsors on $N^\times$
that restrict on every geometric fibre to the standard root cover
is the gerbe of lifts for the pushout extension. Its local objects
are exactly these root covers, and its transition identifications
have the displayed triple obstruction. This identifies its etale
class with the image of the group-extension class. Inverting the
inertia identification or dualizing the
line reverses its sign. Higher formal-coordinate terms have unique
$n$th roots and do not affect the calculation. This argument uses
etale refinements or hypercovers, as in the earlier proof, and does
not assume that a particular Zariski cover computes $H^2$.

If $S$ is an etale $K(\pi,1)$, finite locally constant cohomology
agrees with that of $G$. Apply this at every $\mu_n$ and pass to
the factorial cofinal system. The continuous-cohomology exact
sequence and the bounded Picard-torsion hypothesis verified in
F149-KUM give the claimed continuous degree-two identification.
This explains exactly the extra comparison used in scan~80.

Suppose now that $\operatorname{Pic}(S)$ is finitely generated.
If the extension splits, every finite pushout class is zero, so
every Kummer boundary $\delta_n(\mathcal L)$ is zero. The finite
Kummer sequence then gives $\mathcal L\in n\operatorname{Pic}(S)$
for every positive integer $n$. Their intersection is zero in a
finitely generated abelian group, hence $\mathcal L$ is trivial.
Conversely a conormal basis gives a section of $N^\times\to S$
(the dual normal vector taking value one on that basis). The
induced fundamental-group map splits (R149.3), hence also the
normalizer extension. The root-coordinate calculation of
F149-SPLIT identifies a change of basis with the Kummer action of
its actual global unit. This proves the last assertions.
\end{proof}

This settles the relative comparison in the smooth
characteristic-zero setting of the excerpt, with its explicit
homotopy-exactness hypothesis retained. It does not assert that an
arbitrary fibration has an injective geometric fundamental group,
or replace tame covers by wild covers in positive characteristic.
The intrinsic recognition of the uncompleted geometric tangent
subset in scan~85 remains a different question.
